\section{Composition and Proper Coordination}
\label{sec:composition}

Prior literature presented a puzzle: the ``fix'' for a non-monotone
program is coordination, which is another non-monotone program. A key-value
store with mutable keys is non-monotone~\cite{alvaro2011bloom}. The standard
solution, as in state-machine replication~\cite{schneider1990smr}, it to wrap it 
with a log. Clients send
requests, a log puts the requested commands in one order, and the key-value store
executes the log and answers each request.  The log's specification is also not
monotone: when two commands arrive concurrently, either may be appended next,
and the two resulting logs have no common \outnoun.  
Yet replicated stores are consistent, and they
coordinate only in the log.  Multi-Paxos agrees on each log
entry~\cite{lamport2001paxos}, and the replicas then execute the agreed log---the
key-value store needs no further agreement.

That puzzle is explained by our framework.  Specifications are
tables, so they compose by equijoin, and composition preserves Lifted
Monotonicity.  Because commitments are explicit, we
can also ask whether commitments made at one component are enough to make the
whole composition consistent.  When they are, we call the coordination
\emph{proper} (Section~\ref{sec:proper-coordination}).  Example~\ref{ex:committed-log-composition}
returns to the replicated store.  Proofs are in
Appendix~\ref{app:composition}.
\label{sec:sequential-composition}

Intuitively, sequential composition feeds the output of one specification to
the input of the next.  The composite is again a table: it pairs an input of
the first specification with every output that the second permits for some
intermediate value that the first permits.  Formally, let
$\Spec_1=(I,\InputOrd,M,\preceq_M,S_1)$ and
$\Spec_2=(M,\preceq_M,O,\OutputOrd,S_2)$ be ordered relational specifications
sharing the ordered domain $(M,\preceq_M)$.  Their \emph{sequential
composition} has table
\begin{equation}
  S_2\circ S_1
  =\pi_{I,O}(S_1\Join_M S_2)
  =\{(i,o):\exists m.\ (i,m)\in S_1\land(m,o)\in S_2\},
  \label{eq:spec-composition}
\end{equation}
where $\pi$ and $\Join$ are relational projection and equijoin.  Since $S_1$
is total on $I$ and $S_2$ is total on $M$ (it has a row at every $m\in M$),
every $i\in I$ has a nonempty fiber in $S_2\circ S_1$.

\begin{lemma}[Composition]
\label{lem:composition}
If $\Spec_1$ and $\Spec_2$ satisfy Lifted Monotonicity, then the ordered
specification with table $S_2\circ S_1$ satisfies Lifted Monotonicity.
\end{lemma}

\begin{proof}[Proof sketch]
Take rows of $S_2\circ S_1$ whose inputs have a common \indesc $j$.  The
common-\indesc condition for $\Spec_1$ gives one intermediate value at $j$
that \outrel each row's intermediate value, and that value is a common \indesc
at which the condition for $\Spec_2$ supplies the common \outnoun.  Full proof
in Appendix~\ref{app:composition}.
\end{proof}

By associativity the lemma extends to finite pipelines.  For a function $f$,
its table $\{(i,f(i))\}$ satisfies Lifted Monotonicity exactly when $f$ is
monotone (Section~\ref{sec:lifted-monotonicity}). 

The CRDT observations of Section~\ref{sec:crdt-instance} have this shape: a
client query $q$ over the replica state $s$ is the composition $q\circ s$, and
compositions of monotone functions are monotone
(Corollary~\ref{cor:functional-composition}), so
monotonicity of $q$ on all of the state domain suffices.

\subsection{Proper Coordination}
\label{sec:proper-coordination}

When an upstream component is not monotone, Lemma~\ref{lem:composition} does not apply, and
the composition may need commitments.  Proper coordination asks whether commitments made
upstream are enough on their own.  We describe those commitments by their
effect on the upstream table.  Let $S_1\subseteq I\times M$ and
$S_2\subseteq M\times O$ be as above.  A
\emph{commitment strategy} for $S_1$ (distinct from a participant's strategy,
Definition~\ref{def:system}) fixes at each input $i$ the upstream outputs
committed to at $i$.  Its \emph{committed table} is the set $K\subseteq S_1$
of those rows.  It is \emph{totality-preserving} when $K$ has a row at every
input, so that the commitments never leave an input without a permitted
output.  The resolved end-to-end table is $S_2\circ K$.  Table $K$ is given, not derived: it records which
upstream rows the commitments keep.

\begin{definition}[Proper Coordination]
\label{def:proper-coordination}
A totality-preserving commitment strategy with committed table $K$ is
\emph{proper} for the composition when the resolved specification
$S_2\circ K$ is consistent and every output at every
state of it reachable without commitment is exposable.
\end{definition}

The second clause follows from the first by Complete CALM and the exposure
boundary (Theorem~\ref{thm:complete-calm},
Proposition~\ref{prop:exposure-boundary}): once the upstream commitments are
made, the composition needs no further commitment.  Neither $S_1$, $S_2$, nor
the unresolved composition need be monotone.

\begin{corollary}[Proper Coordination]
\label{thm:proper-coordination}
A totality-preserving commitment strategy is proper for a composition if
and only if the resolved
specification satisfies Lifted Monotonicity.
\end{corollary}

\begin{proof}[Proof sketch]
Complete CALM equates consistency of the resolved specification with Lifted
Monotonicity, and the exposure boundary supplies the second clause.  Full
proof in Appendix~\ref{app:composition}.
\end{proof}

The corollary reduces properness to a test on one table.  A
sufficient condition commits the upstream component to a single output at
each input, chosen so that the outputs grow with the input.

\begin{corollary}[Committing upstream to a monotone function]
\label{cor:monotone-commitment}
If the committed table $K$ of a totality-preserving commitment strategy is the
table of a monotone function and the downstream specification $S_2$
satisfies Lifted Monotonicity, the strategy is proper.
\end{corollary}

\begin{proof}[Proof sketch]
The table of a monotone function satisfies Lifted Monotonicity, so
Lemma~\ref{lem:composition} makes the resolved table monotone.  Full proof in Appendix~\ref{app:composition}.
\end{proof}

The converse fails: a strategy can be proper even when the downstream
specification is not monotone (Appendix~\ref{app:composition}).

\begin{example}[Committed log composed with a key-value store]
\label{ex:committed-log-composition}
Consider again state-machine replication, with every put and get passing
through the log.  Take $I$ the proposal streams received by a single log
authority and $M$ the logs, both ordered by prefix, and $O$ the cumulative
response histories, ordered by extension.  Each command is proposed at most
once.  Let $S_1$ relate each proposal stream to the logs still permitted for
it, and $S_2$ relate a log to the responses obtained by executing it.  Then
$S_1$ is non-monotone, since distinct next entries are incompatible outputs in
one fiber.  So is the key-value store specification $S_2\circ S_1$, since
concurrent commands ordered oppositely give incompatible responses.

Let the authority commit, at each position, to the next command in arrival
order.  Arrival order is part of the input, so $K$ assigns each stream one
log, which extends as the stream does: $K$ is the table of a monotone
function.  Deterministic execution makes $S_2$ monotone.  By
Corollary~\ref{cor:monotone-commitment} the log commitments are proper, and
$S_2\circ K$ is consistent with no commitment in the store.  The binding choice
is the order itself, made upstream.  Multi-Paxos realizes the authority
across replicas~\cite{schneider1990smr,lamport2001paxos}, with the decided
order in place of arrival order.
\end{example}

A downstream component often has input of its own, not only the output of the
component before it.  In a replicated store, for example, not every request
needs to pass through the log.  Some stores serve \emph{bare reads}: a client
sends a get directly to a replica, which answers from the log prefix it has
executed, without going through Multi-Paxos.  Bare reads avoid the cost of
coordination, at the risk of missing a fresher value that has already reached
another replica.  In Appendix~\ref{sec:parameterized-composition} we give
each component a separate input, in an \emph{independently parameterized
composition} for which the results above still hold.  There, bare reads are
consistent when each response records the log position it read from, but not
when the specification demands the latest value
(Example~\ref{ex:bare-reads}).
